\section{The exact recursion may be rooted anywhere}\label{dt:sec-3}

Lower the fixed Frobenius pairing at every vertex of the original closed
source. Every physical turn then identifies equal Newton indices and
retains its factor \(\gamma_e H_N/W_k\). The resulting finite scalar is a
sum over the original edge indices of a product of the original vertex
trace tensors. Choosing a root merely specifies the order of these
finite contractions. It does not move a cut through a tangle, commute
open operators or change a framing.

Consequently the exact recursions of \cite[Theorem~7.1]{Meridian} hold for any chosen root:

\begin{equation}\label{dt:eq-25}
M_{v,N}(u)=\gamma_{e_v}(s)\,
 \mathcal D_{n_v,\sigma_v,N}^{(t)}
 \left[K_N(t,u;s)\prod_{w\in\mathrm{Ch}(v)}M_{w,N}(t)\right]
\end{equation}

at a nonroot vertex, where
\(\mathcal D_{n,\sigma,N}[f]=\phi_\sigma(R_\sigma^n f)\)
uses the actual physical crossing polynomial and its Cartan scalar.
At the root the recursion is

\begin{equation}\label{dt:eq-26}
Z_N=\mathcal D_{n_o,\sigma_o,N}^{(t)}
                  \left[\prod_{w\in\mathrm{Ch}(o)}M_{w,N}(t)\right]
\end{equation}

at the root. The original global closing, framing and reference factors of \cite[Proposition~8.1]{Meridian},
and the original single reduced cut \(H_N\), remain outside \eqref{dt:eq-26}. Edge
gammas are counted once. The exact normal value \(\sigma_v=0\) and the moving
values \(\pm s\) retain their actual orientation types.

Each message is an even trigonometric polynomial of degree at most \(N-1\)
in \(u\), regardless of the subtree. There is no discarded-source error in
\eqref{dt:eq-25}--\eqref{dt:eq-26}.

\begin{figure}[t]
\centering\begin{tikzpicture}[x=1cm,y=1cm,>=Stealth,vertex/.style={circle,draw,fill=white,minimum size=7mm,inner sep=1pt},every node/.style={font=\small}]
\node[vertex] (a) at (-4,1.6) {$w_1$};
\node[vertex] (b) at (-2.2,1.6) {$w_2$};
\node[vertex] (v) at (-3.1,0) {$v$};
\node[vertex] (p) at (-3.1,-1.5) {$p$};
\draw[->] (a)--node[left] {$M_{w_1,N}$} (v);
\draw[->] (b)--node[right] {$M_{w_2,N}$} (v);
\draw[->,thick] (v)--node[right] {$M_{v,N}(u)$} (p);
\node[anchor=west,align=left,text width=7.4cm] at (-.7,.8)
 {$\widehat M_{w,N}(t)=M_{w,N}(t)/M_{w,N}(0)$\\[4pt]
 $M_{v,N}(u)=\gamma_{e_v}\!\left(\prod_wM_{w,N}(0)\right)$\\[2pt]
 $\displaystyle\hspace{7mm}\cdot\mathcal D_v\!\left[K_N(t,u)\prod_w\widehat M_{w,N}(t)\right]$};
\node[draw,rounded corners,align=left,text width=6.8cm,inner sep=7pt] at (2.8,-1.2)
 {At the final root: omit $K_N(t,u)$ and $\gamma_{e_v}$. There is no outgoing polynomial.};
\end{tikzpicture}

\caption{Edges point toward the vertex whose trace is evaluated next.
The exact numbers \(M_{w,N}(0)\) are factored out before applying that
trace. The outgoing polynomial is then divided by its own exact value
at zero. Only a nonroot vertex must supply an outgoing real-decay bound;
the final root evaluates a scalar. This diagram describes contraction
order, not an isotopy of the knot.}\label{dt:recursion-figure}
\end{figure}

\section{The phase class and the variable envelopes}\label{dt:sec-4}

The inductive analytic class is

\begin{equation}\label{dt:eq-27}
h_v(0)=h'_v(0)=0,\qquad h_v(-u)=h_v(u),\qquad
 |h_v''(u)+a_0|\le a_0/20,\quad u\in\mathcal R.
\end{equation}

The leading normalized amplitude \(\alpha_v(u)/\alpha_v(0)\) is holomorphic
and nonzero on a neighborhood of the required smaller compact. Exact
subtree constants are retained:

\begin{equation}\label{dt:eq-28}
\frac{M_{v,N}(u)}{M_{v,N}(0)}
 =\frac{\alpha_v(u)}{\alpha_v(0)}e^{Nh_v(u)}(1+O(N^{-1})),
                     \quad u\in\mathcal R.
\end{equation}

As usual a slightly smaller rectangle can be used for fixed derivative
estimates; the strict margins permit an initially slightly larger one.

The global assertions attached to \(c_v\) are

\begin{equation}\label{dt:eq-29}
\left|\frac{M_{v,N}(t)}{M_{v,N}(0)}\right|
          \le CN^A e^{-Nc_v(1-\cos t)},
                         \quad t\in[-\pi,\pi],
\end{equation}

\begin{equation}\label{dt:eq-30}
\left|\frac{M_{v,N}(t+iy)}{M_{v,N}(0)}\right|
 \le CN^A e^{N[-(c_v/2)(1-\cos t)+y^2]},
                      \quad |y|\le\delta_v.
\end{equation}

For canonical representatives of degree less than \(N\), the real sup,
Chebyshev and full value-and-first-derivative norms differ by polynomial
factors, by \cite[Lemma~2.1]{Summed}. A product of degree a fixed multiple
of \(N\) is first bounded in its real sup norm and then folded in the
correct quotient; no equivalence with finite node data is asserted for
an unreduced higher-degree polynomial. In particular the reduction and
trace retain the repeated-root normals.

Here is the passage from the real bound to the strip. Write the
normalized message as a Laurent polynomial \(P(e^{\ii t})\), of powers
between \(-m\) and \(m\), where \(m<N\). The subharmonic function
\(\log|z^mP(z)|\) on the unit disk is bounded by the harmonic
function \(\log(CN^A)-Nc_v+Nc_v\Re z\). Substitute
\(z=e^{-y+\ii t}\) and divide by \(|z|^m\); for negative \(y\)
use the reciprocal polynomial. Since \(1-e^{-|y|}\le|y|\), this gives

\begin{equation}\label{dt:eq-31}
|M_{v,N}(t+iy)/M_{v,N}(0)|
 \le CN^A e^{N[-c_v(1-\cos t)+(1+c_v)|y|]}.
\end{equation}

For \(|t|\ge r_0\) use \eqref{dt:eq-3}. For \(|t|<r_0\) use the rectangle expansion and

\begin{equation}\label{dt:eq-32}
\Re h_v(t+iy)\le
       -\frac{19a_0}{40}t^2+\frac{21a_0}{40}y^2.
\end{equation}

Both arguments give \eqref{dt:eq-30}, without a conjectured quadratic estimate for
an individual high Newton basis element.

Assume the children satisfy \eqref{dt:eq-27}--\eqref{dt:eq-30}. At a nonroot vertex set

\begin{equation}\label{dt:eq-33}
A_v^{\mathrm{loc}}(t,u)=\psi(t,u)+\sum_wh_w(t),
 \quad
 C_{v,\ell}(u)=\operatorname{stat}_t
       \{A_v^{\mathrm{loc}}(t,u)+ib_vt^2/2-i\ell t\}.
\end{equation}

At the root omit \(\psi\) and \(u\). There are \(d_v\) incidences in \(t\). Hence

\begin{equation}\label{dt:eq-34}
A_v^{\mathrm{loc}}(t,u)=\psi(0,u)-a_vt^2/2+E_v(t,u),
       \quad |E_{v,tt}|,\ |E_{v,t}/t|\le a_v/20.
\end{equation}

The root constant in this formula is zero. For \(0\le\ell\le L_v=\lfloor t_0|D_v|\rfloor\)
the equation \(D_vt=-i\ell+E_{v,t}\) contracts in \(|t|\le b_0\). Indeed
\(t_0+b_0/20<b_0\). The actual solutions lie in the smaller disk
\(|t|\le t_0/(19/20)\), leaving room for the translated contours.

Writing \(E_{v,t}=tQ_v\), with \(|Q_v|\le a_v/20\), gives
\(t_{v,\ell}=-\ii\ell/(D_v-Q_v)\). Its imaginary part is
\(-\ell\Re(D_v-Q_v)/|D_v-Q_v|^2\). The numerator lies
between \(19a_v/20\) and \(21a_v/20\), and the denominator
between \((19|D_v|/20)^2\) and \((21|D_v|/20)^2\).
Therefore

\begin{equation}\label{dt:eq-35}
\frac{380}{441}\frac{\ell a_v}{|D_v|^2}
 \le-\Im t_{v,\ell}\le
 \frac{420}{361}\frac{\ell a_v}{|D_v|^2},
\end{equation}

also for complex parent \(u\). No child phase is asserted to be real on
the real axis. Its nonlinear derivative is \(t\) times a function bounded
by \(a_v/20\), so the real part of the denominator remains positive.

The normalized phase \(h_v=C_{v,1}(u)-C_{v,1}(0)\) satisfies

\begin{equation}\label{dt:eq-36}
h'_v=\psi_u(t_{v,1},u),\qquad
 h''_v=\psi_{uu}
       -\frac{\psi_{ut}^2}{A_{v,tt}^{\mathrm{loc}}+ib_v}.
\end{equation}

Using \eqref{dt:eq-17} and \(|A_{v,tt}^{\mathrm{loc}}+ib_v|\ge(19/20)|D_v|\ge(19/20)a_0\) proves \eqref{dt:eq-27},
since \(1/50+20/(19\cdot2500)<1/20\). This is uniform on the rectangle.
Charge ordering and the current-vertex drop are

\begin{equation}\label{dt:eq-37}
\Re C_{v,\ell}(u)\le\Re C_{v,1}(u)
  -\frac{190}{441}\frac{a_v}{|D_v|^2}(\ell^2-1),
 \quad
 0\le\Re\psi(0,u)-\Re C_{v,1}(u)\le\beta_v.
\end{equation}

At the root again replace \(\psi(0,u)\) by zero.

\section{Actual trace tails with the directed budgets}\label{dt:sec-5}

Divide every child by its exact nonzero value \(M_{w,N}(0)\). The induction
uses those complex constants exactly; it never estimates them by their
complete-source magnitudes. Equations \eqref{dt:eq-15} and \eqref{dt:eq-30} imply

\begin{equation}\label{dt:eq-38}
|\text{normalized trace input}(t+iy,u)|
 \le CN^A
   e^{N[\Re\psi(0,u)-A_vt^2+B_vy^2]},
                      \quad |y|\le\delta_v^{\mathrm{in}}.
\end{equation}

The root constant is zero. The factor \(2/\pi^2\) in \eqref{dt:eq-4} is exactly the
inequality \(1-\cos t\ge2t^2/\pi^2\) on the original interval.

Write \(\eta=(-1)^d\). For the actual input, an even trigonometric
polynomial of any fixed-multiple degree, the exact trace is
\begin{equation}\label{dt:trace-grid}
 \mathcal D_{n,\sigma,N}[f]
 =\frac{\eta N C_{N,n}(\sigma)}{\sin(2\pi\sigma)}
   \sum_{j=0}^{N-1}\sin\tau_j\,f(\tau_j)
                      e^{\ii nN\tau_j^2/(4\pi)},
 \quad\begin{cases}
 \tau_j=\dfrac{2\pi}{N}(d/2+\sigma-j),\\
 C_{N,n}(\sigma)=q_N^{n(d^2/4-\sigma^2/2)}.
 \end{cases}
\end{equation}
This is \cite[equation~(6.4)]{Meridian}, initially away from colliding
grids; its holomorphic limit includes every normal. To recall the exact
Fourier conversion, use
\[
 \mathcal K_\sigma(t)=
 \frac{\ii N e^{\ii Nt}}{e^{\ii Nt}-\eta e^{2\pi\ii\sigma}}
\]
on the rectangle with sides \(\pm\pi\pm\ii y_v+2\pi\sigma/N\).
The kernel has residue one at the \(N\) enclosed grid points and equals
\(\ii N e^{-Nz}/(1+e^{-Nz})\) on a vertical side
\(t=\pm\pi+\ii z+2\pi\sigma/N\). Its upper and lower horizontal
series are
\[
 -\ii N\sum_{\ell\ge1}\eta^\ell
       e^{-2\pi\ii\ell\sigma}e^{\ii\ell Nt},\qquad
 \ii N\left(1+\sum_{\ell\ge1}\eta^\ell
       e^{2\pi\ii\ell\sigma}e^{-\ii\ell Nt}\right).
\]
Move each entire term to a fixed horizontal segment, retaining its short
endpoint connectors. Reflect the upper integral by \(t\mapsto-t\),
with the boundary orientation; \(\sin t\,f(t)e^{\ii nNt^2/(4\pi)}\)
is odd. The paired coefficients are therefore
\begin{equation}\label{dt:eq-39}
 c_{\ell,N}(\sigma)=\eta^{\ell-1}
 U_{\ell-1}(\cos2\pi\sigma),\qquad c_{1,N}=1,
\end{equation}
with common prefactor \(-N^2C_{N,n}/(\pi\ii)\).
The zero-mode integral on the real interval vanishes by oddness, but its
endpoint currents remain. These algebraic steps are also written in
\cite[Section~6.1]{Summed}; they impose no condition on the twists or
sum over vertices. In particular there is no later sum of unexamined
equal-action signs.

For \(\ell>L_v\), moving the single-grid rectangle by \(y_v\) costs, after
maximizing the chirp term in the real variable, at most \(C_v y_v^2\).
Since \(\ell>L_v\) implies \(\ell>t_0|D_v|\), \eqref{dt:eq-6} gives a tail loss at least

\begin{equation}\label{dt:eq-40}
t_0|D_v|y_v-C_vy_v^2\ge\frac12t_0|D_v|y_v.
\end{equation}

Successive charges have the additional factor \(\exp(-Ny_v)\). The exact
fixed-compact growth of \eqref{dt:eq-39} is absorbed for all sufficiently large N.
This starting level may depend on the fixed vertex.

At an original endpoint \(\pm\pi\), the chirp and finite-charge connectors cost
at most

\begin{equation}\label{dt:eq-41}
(\pi|b_v|+t_0|D_v|)y_v+B_vy_v^2
                            \le A_v\pi^2/4.
\end{equation}

Thus their loss is at least \(A_v\pi^2/2\). This includes the original
zero-mode current. Nothing is inferred by dropping the endpoint terms
of the nonperiodic chirp.

For a retained charge, restrict its real integral to \([-r_0,r_0]\).
Its complement loses \(A_vr_0^2\) by \eqref{dt:eq-38}. Translate the retained interval
to \(\Im t_{v,\ell}\). Its endpoints remain in the analytic rectangle. The
original local Young inequality, now on this larger real interval,
is

\begin{equation}\label{dt:eq-42}
\begin{aligned}
 \Re\{A_v^{\mathrm{loc}}(t-iz,u)+ib_v(t-iz)^2/2-i\ell(t-iz)\}
 &\le\Re\psi(0,u)-a_vt^2/8\\
 &\quad+\left[\frac{21}{40}a_v+
                 \frac{2b_v^2}{(14/5)a_v}\right]z^2-\ell z.
\end{aligned}
\end{equation}

For \(0\le z\le-\Im t_{v,\ell}\), the last line after its quadratic term is
at most \(-\ell z/8\), since \((420/361)\max(21/40,5/7)=300/361<7/8\).
Every such connector therefore loses at least \(a_vr_0^2/8\).
The translated horizontal saddle has negative real quadratic part
at most \(-(19/20)a_v\). The physical sine factor is nonzero at \(t_{v,\ell}\).

Equations \eqref{dt:eq-40}--\eqref{dt:eq-42} give an error of at most
\(CN^A\exp N[\Re\psi(0,u)-\eta_v]\). Nonminimal charges are their saddle
terms plus this absolute error; no relative approximation is asserted
at a much lower nonminimal action. Conditions \eqref{dt:eq-8} and \eqref{dt:eq-37} put the
whole conditional error strictly below the first charge. Equation
\eqref{dt:eq-39} gives its nonzero coefficient for both parities. This establishes
\eqref{dt:eq-28} and nonvanishing of the exact \(M_{v,N}(0)\).

For the child inputs, their exact analytic logarithms differ from the
limiting phases by \(O(N^{-1})\), with fixed derivatives on the required
smaller compact. The finite product of those corrections is retained
in the amplitude, or moves the saddle by \(O(N^{-1})\). As the tree is fixed,
it spends no new exponential budget. The leading amplitude is a product
of nonzero edge, child and Gaussian factors. All N-dependent physical
crossing and carrier phases remain in \(M_{v,N}(0)\); its normalized leading
amplitude does not conceal such phases.

The full physical trace and the first-charge approximation are
holomorphic in the formal meridian variables. As in the original
proof, first take parameter circles avoiding collisions and then apply
the maximum principle to their whole difference. This covers every
actual \(\sigma_v=0,\pm s\), including normal limits. No isolated divided
current is declared removable without its companions.

\section{Regenerating the outgoing real bound}\label{dt:sec-6}

Lemma~2.1 of \cite{Summed} compares the real sup norm, the
Chebyshev norm, and the complete value-and-first-derivative norm by
polynomial factors. It also bounds the physical fixed-power crossing
and the full quotient trace. These estimates apply to polynomial
products of degree a fixed multiple of \(N\); that multiple can depend
on the degree of this fixed vertex. For real \(u\), \eqref{dt:eq-14},
the polynomial real sup bounds of the normalized children, the actual
fixed-power crossing polynomial, Chebyshev folding in the correct
quotient, and the full top-coefficient functional give

\begin{equation}\label{dt:eq-43}
|M_{v,N}(u)|\le CN^A
       \left|\prod_wM_{w,N}(0)\right|
                  e^{N[2G_0-k_0(1-\cos u)]}.
\end{equation}

This retains repeated-root derivatives and all fixed normals. Dividing
by the nonzero first-charge value at zero costs at most \(\beta_v\). For
\(|u|\ge r_0\), \eqref{dt:eq-9} absorbs that cost and yields \eqref{dt:eq-29}. For \(|u|<r_0\), \eqref{dt:eq-28}
and \eqref{dt:eq-32} give the stronger estimate, since \(c_v<k_0\ll19a_0/20\).
Equation \eqref{dt:eq-30} follows by Section~\ref{dt:sec-4}. Thus the induction closes with
the selected directed budgets, rather than a common decay constant.

There is no outgoing real bound at the final root. Imposing \eqref{dt:eq-9} there
would be an unnecessary restriction. Its only normalization comparison
is the conditional trace estimate \eqref{dt:eq-8}.

\section{Original action, geometry, and scope}\label{dt:sec-7}

Restore the exact subtree constants successively. Stationary
elimination in \eqref{dt:eq-33} gives the original joint equations

\begin{equation}\label{dt:eq-44}
\sum_{e\sim v}\psi_{e,\tau_v}
                      +\frac{in_v\tau_v}{2\pi}=i.
\end{equation}

The exact child constants telescope to the original critical value

\begin{equation}\label{dt:eq-45}
\mathcal C(\mathbf1)=\operatorname{stat}_{\boldsymbol\tau}
 \left[\sum_e\psi_e+
       \frac{i}{4\pi}\sum_vn_v\tau_v^2-i\sum_v\tau_v\right].
\end{equation}

Back-substitute from the final root: its first-charge angle belongs to
the rectangle, so each child's selected parent value does also, and
the same statement continues down the finite tree. Every restored angle
has modulus at most \(1/[(19/20)|D_v|]\). Thus the vector satisfying
\eqref{dt:eq-44} lies inside the full simultaneous small-angular uniqueness
polydisk, not just the conditional disks used during elimination.
The existing contraction identifies it with the same branch as the
marked positive charge path of \cite[Theorem~1.1]{Geometry}. In fact
\(20/(19|D_v|)<10^{-2}\). Indeed
\eqref{dt:eq-8} implies

\begin{equation}\label{dt:eq-46}
\sqrt{d_v^2+n_v^2}=2\pi|D_v|\ge 2\pi\,2^{15}>630.
\end{equation}

\cite[Theorem~1.1]{Geometry} applies with no sum condition.
The edge functions agree: their derivatives are \eqref{dt:psi} and
their complete values satisfy \(2G_0=v_{\rm oct}/\pi\), as in
\cite[Section~7]{Summed}. Write \(n_v=\epsilon_v-2p_v\), with
\(\epsilon_v\in\{-1,0,1\}\), in the original disk frame. The
primitive periods supplied by the geometric theorem are
\[
 U_v=-2\pi\Psi_{\tau_v}-\ii\epsilon_v\tau_v,
 \qquad V_v=2\ii\tau_v.
\]
Equation \eqref{dt:eq-44} therefore gives the full continued relation
\(U_v+p_vV_v=-2\pi\ii\). The imported theorem proves its primitive
winding, smooth metric completion, and complete main cusp; we do not
infer these facts from an exponentiated equation.
Its original ordered two-octahedron cut, odd seam and calibrated action identify \(2\pi\Re\mathcal C(\mathbf1)\) with the volume of the actual
filled knot complement. This is a geometric input, not a deduction
from the quantum saddle alone.

There are \(|E|\) edge Gaussians \(N^{-3/2}\) and \(|V|\) trace Gaussians \(N^{3/2}\);
\(|V|=|E|+1\) leaves the original \(N^{3/2}\) power with nonzero amplitude.
The exact first Fourier coefficient is one. The retained edge insertion
has leading value one, and the scalar ledger remains precisely \eqref{dt:scalars}.
At \(s=0\), all edge gammas and paired insertions equal one.
The exact outside scalar is
\[
 \Omega_N=\theta_d^{-w(D)}\prod_vC_{N,n_v}(0)
       =q_N^{d^2(2w(D)+\sum_v n_v)/4},
 \qquad \theta_d=q_N^{-d^2/2}.
\]
The remaining leading coefficient is independent of \(N\), nonzero,
and the same through both parities. This proves the full leading
asymptotic in Theorem~\ref{dt:main}.
For bounded \(s\), the first-charge coefficient is still one, the
paired insertion has leading value one, and all normalized leading
amplitudes are unchanged. Restoring the exact scalar ledger of
\cite[Proposition~8.1]{Meridian} gives \(A_N(s)H_N(s)\).
The finite product of relative \(1+O_R(N^{-1})\) factors proves
\eqref{dt:profile}. Its nonvanishing follows uniformly on every
fixed disk \(R<1\), since \(A_N\) and \(H_N\) have no zeros there.
Together with the geometric identification above this proves
Theorem~\ref{dt:main}.

The criterion allows different strip and real-decay budgets and a
different choice of final vertex. It is not the assertion that arbitrary
capped short-twist islands can be substituted as black-box messages.
Such an island requires a multi-boundary version of the conditional
normalization proved here for an actual edge. The pointwise source
estimates on the older tiny disk alone do not supply that extension.
